%======================================================================
\section{讨论}
\label{sec:discussion}
%======================================================================

\paragraph{羊群效应是病态吸引子，而非措辞缺陷。}
研究一与跨模型复现收敛于一幅结构性的图景：多轮 LLM 评审团并非仅仅收敛——它们收敛于模型特定的吸引子。对 DeepSeek-V3 与 R1 而言，吸引子是拒稿；对 GLM-4-Flash 而言，是带着骑墙语言的接收。吸引子的引力强大到一项显式的、方向中立的调控协议（研究 2b）几乎改变不了任何事：被指示除非遭到具体驳倒否则坚持的少数派，仍与未调控对照以相同的速率投降。这呼应了问答领域的发现——静态反谄媚必须作用于机制\emph{参数}（EV 权重抑制）而非提示——并将其推广：当底层模型先验是从众的时，连行为提示也会失效。因此，对 LLM 评审团的有效调控必须作用于聚合层或审议结构层，而非修辞层。

\paragraph{校准优先原则。}
最清晰的正面结果同时也是成本最低的：拒稿偏置——单智能体先验的扭曲——在聚合时即被修复。Stackelberg 优化的 $\theta^*$ 自主发现了这一点，在共识分数上放置了 $+1.04$ 的校准项，而更简单的重设阈值均值捕获了几乎相同的一致率。实践原则：在增加审议之前，先校准聚合。未经校准的审议会通过棘轮放大先验的偏置；经校准的独立聚合在我们测试的所有讨论协议中，单位调用的决策一致率已然最优。

\paragraph{讨论作为按需资源。}
如今在两个领域出现了同一模式：大多数实例不需要讨论（70\% 的评审团在第一轮后即一致；在问答基准上，大多数问题对拓扑不敏感），且无指导的讨论可能主动造成损害。路由将这一观察转化为成本节约：评审领域的路由器仅在 17\% 的论文上支出审议预算，并匹配最优固定协议。10 个百分点的预言机差距标志着前沿——更好的前瞻特征（论文难度信号、超越熵的评审分歧语义）可能将其解锁。

\paragraph{行为健康作为主要终点。}
本文的一项方法论贡献是终点原则。决策一致率在任何条件或干预下几乎纹丝不动（全部落在 59--65\% 区间），而行为指标的变动幅度超过 50 个百分点且 $p{<}10^{-28}$。倘若我们只测量决策，研究一将报告"讨论无害"——这恰是错误的结论。任何对多智能体判断系统的评估都应报告少数派存活、方差保持与漂移，并将准确率降级为安全约束。当真值含噪声时（如官方评审决议），这一点尤为重要。

\paragraph{向同行评审之外迁移。}
三个层次可映射到通用的多 LLM 判断管道。校准聚合适用于任何多信号融合为决策的场景——推荐系统中的排序集成、多评审官评估、红队投票——任何共享先验偏置需要学习校正项的地方。行为健康指标（消灭率、方差坍缩、漂移、翻转不对称）构成了多智能体管道的\emph{审议健康监视器}：它们仅从记录的立场即可计算，可作为"先审议后行动"的智能体系统的生产级可观测信号。协议路由可推广到算力分配：哪些查询值得多模型审议，哪些应由单次经校准的调用回答。我们尚未实证验证这些迁移；它们是自然的后续研究。

\paragraph{与问答领域验证的关系。}
集体判断有两个问题——\emph{决策对不对}与\emph{过程健不健康}——框架的两次验证各回答一个。在问答基准上，同一框架验证了路由层（87.6\% 准确率，比最佳固定拓扑高 +4.5 个百分点）与 AEIC 的结构化分解前提（显式结论字段带来 +28 个百分点），而其质量信号被证明衡量的是连贯性而非正确性。评审领域改变了问题：在含噪真值与行为风险之下，框架的价值从准确率增益转向诊断（使羊群效应可测量）、校准（修复先验偏置）与成本感知的编排。两次验证互为补充：问答表明框架能优化可优化之物；评审表明框架能调控须保护之物。

\paragraph{局限性。}
(1)~评审团按模型同质；异质评审团（混合模型、混合人设）的羊群行为可能不同——我们的三模型研究采样了后端，但未在评审团内部混合。
(2)~官方决议是含噪真值；我们通过安全终点原则与分层报告缓解，但一致率的绝对值应读作"与人类评审团的一致"，而非正确性。
(3)~GLM 的立场 94\% 由评分推导（骑墙结论），故其行为指标部分反映评分动态。
(4)~研究 2b 只测试了一种特定的提示层协议；更强的机制级调控（如聚合时少数派保护、结构化驳倒轮）仍未检验。
(5)~路由器的特征是第一轮统计量；前瞻难度估计较浅，预言机差距提示存在可观的上行空间。
(6)~全部数据来自单一会议单一年份（ICLR~2024）；跨会议的普遍性尚待验证。
